\subsubsection{A non-private $L_p$ sampling protocol with $\tilde{O}(1)$
communication} \label{sec:insecureLp}

\floatname{algorithm}{Protocol}
\begin{algorithm}[htbp]

\paragraph{Inputs:} Parties $P_1$ and $P_2$ have inputs $\w_1$ and $\w_2$
respectively. 

\medskip

The protocol proceeds as follows:
\BEN
    \item Parties run $\Pi_{\mathsf{ignore}}$ 
    with inputs $\vec{w}_1, \vec{w}_2$
    that samples from $D_{\mathsf{ignore}, p}(\vec{w}_1, \vec{w}_2)$ and obtain output $i$.
    
    \item For $b \in \{1, 2\}$, $P_b$ sends $w_{b,i}, ||\vec{w}_b||_p^p$.
    Both parties compute 
    \[
    \Pr_{D_{\mathsf{ignore}, p}(\vec{w}_1, \vec{w}_2)}[i] = \frac{w_{1,i}^p + w_{2,i}^p}{||\vec{w}_1||_p^p + ||\vec{w}_2||_p^p}
    \quad
    \mbox{and}
    \quad
    f_c(\vec{w}_1, \vec{w}_2, i) = \frac{(w_{1,i} + w_{2,i})^p}{||\vec{w}_1||_p^p + ||\vec{w}_2||_p^p}
    \]
    \item Parties output $i$ with probability 
    \begin{align*}
        \frac{f_c(\vec{w}_1, \vec{w}_2, i)}{2^{p-1} \cdot \Pr_{D_{\mathsf{ignore}, p}(\vec{w}_1, \vec{w}_2)}[i]} &= \frac{c \cdot \Pr_{D_{L_2}(\vec{w}_1, \vec{w}_2)}[i]}{2^{p-1} \cdot \Pr_{D_{\mathsf{ignore}, p}(\vec{w}_1, \vec{w}_2)}[i]}\\
        &= \frac{\Pr_{D_{L_2}(\vec{w}_1, \vec{w}_2)}[i]}{2^{p-1}/c \cdot \Pr_{D_{\mathsf{ignore}, p}(\vec{w}_1, \vec{w}_2)}[i]}
        \end{align*}
    and otherwise return to step 1.
    
\EEN
\caption{Protocol for exact $L_p$ sampling ($\Pi_{L_p}$)}
\label{fig:lpexact}
\end{algorithm}
In this section we present a $\tilde{O}(1)$ sampling 
protocol for $L_p$ sampling for constant $p$.
We present only the insecure version, extending it
to a secure sampling protocol can be done
entirely analogously to the construction for $L_2$ sampling
given in Section~\ref{sec:secure}.

Given input vectors $\vec{w}_1, \vec{w}_2$,
$L_p$ sampling refers to sampling from the distribution
$D_{L_p}(\vec{w}_1, \vec{w}_2)$ with the following
probability mass function:
$$\Pr[ D_{L_p}(\vec{w}_1, \vec{w}_2) \mbox{ samples } i ] 
= \frac{(w_{1,i} + w_{2,i})^p}{ \sum_j (w_{1,j} + w_{2,j})^p }
= \frac{(w_{1,i} + w_{2,i})^p}{ \| \w_1 + \w_2 \|_p^p }. $$

We begin by defining and showing how to sample from a helper
distribution $D_{\mathsf{ignore}, p}$.

\begin{definition}
For input vectors $\vec{w}_1, \vec{w}_2$,
let $D_{\mathsf{ignore}, p}(\vec{w}_1, \vec{w}_2)$ be the distribution that ``ignores'' the cross term in $D_{L_p}(\vec{w}_1, \vec{w}_2)$.
I.e.~$D_{\mathsf{ignore}, p}(\vec{w}_1, \vec{w}_2)$
samples index $i \in [n]$
with probability $\frac{w^p_{1,i} + w^p_{2,i}}{||\vec{w}_1||_p^p + ||\vec{w}_2|_p^p}$.
\end{definition}

\begin{lemma}\label{lem:ignore_2}
There exists a protocol $\Pi_{\mathsf{ignore}}$ for sampling from
$D_{\mathsf{ignore}, p}(\vec{w}_1, \vec{w}_2)$ with $\tilde{O}(1)$
communication.
\end{lemma}

\begin{proof}
Let $\vec{w'}_b = (w_{b,1}^p, \ldots, w_{b,n}^p).$
The lemma follows by observing the following:
$$D_\ig(\w_1, \w_2) = D_{L_1}(\vec{w'}_1, \vec{w'}_2).$$ \qed
%Each party samples from its own distribution (i.e., $P_b$ outputs $i$
%with probability $\frac{w_{b,i}^2}{||\vec{w}_b||^2}$).  Then the
%parties do a weighted coin flip to choose which of these two values is
%output ($P_b$'s value is chosen with probability
%$\frac{||\vec{w}_b||_2^2}{||\vec{w}_1||_2^2+||\vec{w}_2||_2^2}$.)  It
%is easy to see that this protocol only requires $\tilde{O}(1)$
%communication\anote{Should we call this logarithmic?} and produces the
%correct output distribution.
\end{proof}

\begin{definition}
For $i \in [n]$, let the corrective parameter function be defined as 
$$f_c(\vec{w}_1, \vec{w}_2, i) := \frac{(w_{1,i} +
w_{2,i})^p}{||\vec{w}_1||_p^p + ||\vec{w}_2||_p^p}.$$
\end{definition}


\begin{definition} 
The constant 
$c := c(\vec{w}_1, \vec{w}_2)$ 
is defined as
\[
c(\vec{w}_1, \vec{w}_2) := \frac{||\vec{w}_1 + \vec{w}_2||_p^p}{||\vec{w}_1||_p^p + ||\vec{w}_2||_p^p}
\]
This ensures that 
for every $i$,
$f_c(\vec{w}_1, \vec{w}_2, i) = c \cdot \Pr_{D_{L_p}(\vec{w}_1, \vec{w}_2)}[i]$.
\end{definition}

\iffalse
\begin{proof}
Define 
$c(\vec{w}_1, \vec{w}_2) := \frac{||\vec{w}_1 + \vec{w}_2||_2^2}{||\vec{w}_1||_2^2 + ||\vec{w}_2||_2^2}$.
Clearly, by definition of $f_c$ and $D_{L_2}(\vec{w}_1, \vec{w}_2)$ it is the case that
$f_c(\vec{w}_1, \vec{w}_2, i) = c(\vec{w}_1, \vec{w}_2) \cdot \Pr_{D_{L_2}(\vec{w}_1, \vec{w}_2)}[i]$.
It remains to show that $c(\vec{w}_1, \vec{w}_2) \in [1,2]$.
For the upper bound,
\begin{align*}
    c(\vec{w}_1, \vec{w}_2) &= \frac{||\vec{w}_1||^2_2 + 2\langle \vec{w}_{1}, \vec{w}_{2} \rangle + ||\vec{w}_2||_2^2}{||\vec{w}_1||_2^2 + ||\vec{w}_2||_2^2}\\
    &\leq \frac{2(||\vec{w}_1||^2_2 + ||\vec{w}_2||_2^2)}{||\vec{w}_1||_2^2 + ||\vec{w}_2||_2^2}\\
    &= 2.
\end{align*}
For the lower bound,  $c(\vec{w}_1, \vec{w}_2) \geq 1$, since $2\langle \vec{w}_1, \vec{w}_2 \rangle \geq 0$ (since in our setting $\vec{w}_1, \vec{w}_2$ are positive).
\end{proof}
\fi

The following lemma will be useful for arguing the validity of the
final protocol.

\begin{lemma} \label{lem:close_dist_2}
For all $i \in \mathsf{supp}(D_{L_p}(\vec{w}_1, \vec{w}_2))$, 
$$\Pr_{D_{L_p}(\vec{w}_1, \vec{w}_2)}[i] \leq 2^{p-1}/c \cdot \Pr_{D_{\mathsf{ignore}, p}(\vec{w}_1, \vec{w}_2)}[i].$$
\end{lemma}

The proof is found in Appendix~\ref{apx:pfLemma5}.

We now present the $L_p$ sampling protocol $\Pi_{L_p}$ in Protocol~\ref{fig:lpexact}. We show the correctness and efficiency of the protocol below. 


\begin{lemma} \label{lem:lp}
With all but negligible probability, on inputs $\vec{w}_1$ and $\vec{w}_2$, protocol $\Pi_{L_p}$ samples exactly correctly from $D_{L_p}(\vec{w}_1, \vec{w}_2)$.
Further, for any constant $p$, the protocol has communication
$\tilde{O}(1)$.
\end{lemma}

The proof is found in Appendix~\ref{apx:pfLemma6}.
We note that this result strictly generalizes
Lemma~\ref{lem:insecure_L2_prot}.
In particular, setting $p=2$ in the above
protocol yields a protocol with exactly the 
same parameters as the $L_2$ sampling protocol.


\iffalse
\begin{remark}
Note that the protocol and analysis above
did not require that vectors $\vec{w}_1, \vec{w}_2$ are normalized.
I.e.~we do not require that $||\vec{w}_1||_1$
or $||\vec{w}_2||_1$ are equal to $1$
or to each other.
\end{remark}
\fi